\section{Experiments}

The character-task control--value development Block reduces the model
from 41,593 to 20,721 parameters under the same 300-update all-token
protocol. Its paired 64/128/256-context perplexities are
10.593/10.498/10.560, versus 10.358/10.251/10.253 for the earlier Block.
Training-forward work decreases from 21.715 to 10.342 billion MAC proxies,
but measured training wall time changes from 75.96 to 79.29 seconds.
The GRU and official Mamba2 still obtain lower perplexity and lower
recorded training time. These distinct model versions are development
comparisons, not pooled Primary results or a complete-cost victory.
\label{sec:experiments}


A later three-arm memory-interface development test completed 7,200 new
updates. The 10,351-parameter joint repair achieved 62.70\% new-task IID,
but only 5 of 192 label-changing query pairs were correct at both endpoints.
Its 26.916\,s full-stream update time exceeded the original GRU32's
17.281\,s despite fewer parameters. A frozen VALID content-direction
intervention with exactly matched write weights and coverage reduced
62.50\% accuracy to 41.99\%; strength-matched cached-direction
compensation obtained 39.06\%. These measurements motivate continued
binding research, not a Primary acceptance or comparative efficiency claim;
complete protocols and costs are retained in Appendix~\ref{sec:prometheus-block}.

Our evaluation asks whether contextual computational primitives improve
acquisition, retention, or resource use relative to strong conventional
architectures. Topological components receive a separate causal test when
included. The primary experimental object is the latest fully fused
Prometheus architecture after general Block validation. Monodromy is an
optional component selected by measured benefit; its independent advantage
is not a prerequisite for the Primary freeze. Its core mechanism,
code version, and evaluation configuration will be frozen before retraining
across all five evaluation layers. No publication results for the to-be-frozen
Primary are available at the present manuscript checkpoint. Earlier coordinate-renderer
streams and small integrated diagnostics are development evidence in
Appendices~\ref{app:benchmarking} and~\ref{app:evidence}; they are not pooled
with Primary results or used to supply its performance or efficiency claims.

\subsection{Observed development results}

We first tested the same continuous core on three-stage contextual programs,
sampled character prediction, and initial real-input recognition. These are
fixed-budget, single-seed development measurements, not formal Primary results.
The context stream exposes raw tokens and an initial state, with terminal labels
but no task identifier, correct route, or intermediate state. Training lengths
are 2--12; separate tests use lengths 24--48 and structured unseen compositions.
Each phase uses 600 updates of 64 examples. Actual phase-0 weights, Adam state,
and RNG were restored for later phases; phase0 was not retrained.

\begin{table}[t]
\centering\small
\begin{tabular}{lrrrrr}
\toprule
Model & Parameters & Task0 & Task1 & Task2 & Train s\\
\midrule
Continuous dynamic &28,365&82.81&90.23&91.99&92.21\\
Continuous static caller &28,365&81.45&94.92&94.14&79.12\\
GRU44 &25,572&93.95&95.90&93.36&10.25\\
Transformer32 &26,728&48.63&79.30&83.01&13.76\\
\bottomrule
\end{tabular}
\caption{Final IID accuracy (\%) after three development phases, without
replay. Train s sums actual gradient-step wall time, not total work or exclusive
throughput. The static caller retains contextual primitives and memory.}
\label{tab:development-context}
\end{table}

The dynamic model forgets part of task0 and does not uniquely outperform the
static caller. Distribution checks are essential: a TRAIN-majority constant
predictor scores 69.73\% on task2's long test, equal to the dynamic model's
69.73\%. Its structured-unseen scores are 25.78\% and 25.98\%, respectively.
Thus high IID accuracy alone does not establish general memory computation.

All-token character development uses 300 updates of 32 length64 windows,
614,400 next-character labels per method. A causal-mask Transformer and GRU
receive identical windows and batches. Future-token perturbations test actual
causality before training. The 512 frozen TEST endpoints are paired across
lengths64/128/256; length comparisons score the same final 64 labels. These are
sampled-window character PPLs, not full-corpus or word-level PPLs.

\begin{table}[t]
\centering\small
\begin{tabular}{lrrrrr}
\toprule
Model &P&PPL64&Train s&Driver s&Forward G MAC\\
\midrule
Dynamic &41,593&10.358&75.96&119.47&21.715\\
Static &41,593&10.306&75.24&118.74&21.715\\
Detached &41,593&10.243&74.07&114.41&21.715\\
Transformer36 &39,664&10.533&6.75&12.26&29.816\\
GRU44 &32,660&8.664&2.49&7.96&16.977\\
Official Mamba2 &41,116&7.173&3.56&4.58&---\\
Control--value Block &20,721&10.593&79.29&117.67&10.342\\
Same Block, training Graph &20,721&10.593&7.62&47.72&10.342\\
GRU34 &21,180&9.449&1.09&5.48&10.612\\
Transformer24 &19,564&11.414&3.12&7.78&15.630\\
\bottomrule
\end{tabular}
\caption{Actual all-token language development. Driver time includes measured
evaluation/checkpoint work but excludes process startup. MAC proxies count
training forwards only, not backward or total FLOPs; all raw costs are retained.}
\label{tab:development-ar-cost}
\end{table}

The two smaller conventional widths were fixed prospectively from parameter
formulas, with no accuracy gate or TEST-based configuration selection.
They use the original windows, feedback, Adam, clipping and 300-update budget.
This parameter-matched comparison narrows the apparent weight advantage:
the 20,721-parameter Block remains worse than GRU34 in paired perplexity and
recorded training time. All widths remain separate development runs;
they are not pooled as seeds or substituted for a frozen Primary.
The later Graph run retains the exact core, data, labels and 300 batch indices.
Actual forward/backward replay with eager Adam reduces update wall time by
10.40 times, while preserving paired perplexity to the reported precision.
Setup, diagnostic gradients, memory pools and evaluation are charged separately:
32.16 seconds of evaluation still dominates its 47.72-second driver.
This is an execution improvement, not a new learning rule or a baseline victory.

With contexts of 64, 128, and 256 characters, dynamic common-suffix PPL is
10.358, 10.251, and 10.253; Transformer yields 10.533, 15.185, and 14.736,
and GRU yields 8.664, 8.552, and 8.552. Stability under longer
inputs alone does not prove useful distant-memory use. The earlier terminal-prefix
protocol and serial eager/hard-call inference audits remain separately reported
in the appendix and are not pooled into this new protocol. They show no
hard-call speedup. Cached streaming serving has not been established.
The subsequently executed official Mamba2 baseline uses the same 300-update
protocol and achieves common-suffix PPL 7.173/7.048/7.049 with 41,116 parameters.
Its 3.56-second training and 4.58-second driver times exclude 82.67 seconds of
measured cold compilation/priming; the entire coordinator launch takes 93.31
seconds and creates 55.19 MB of kernel cache. Its exact selective-SSD and
backward FLOPs remain unmeasured. Earlier arms have their own recorded
concurrency, so this is not an isolated-throughput ratio. These results
provide no broad competitive advantage for the current Block.

\begin{table}[t]
\centering\small
\begin{tabular}{lrrrr}
\toprule
Model &Speech acc.&Image acc.&Speech/image P&Train s (S/I)\\
\midrule
Dynamic &29.77&42.10&32,280/40,793&97.51/66.28\\
Static caller &71.17&40.10&32,280/40,793&98.49/66.05\\
Detached &69.38&39.00&32,280/40,793&97.84/65.97\\
GRU64 &76.25&44.60&54,819/59,556&1.60/1.92\\
Transformer64 &76.48&49.90&105,379/110,116&7.92/7.85\\
CNN &59.77&44.30&35,683/38,852&.96/1.05\\
\bottomrule
\end{tabular}
\caption{Initial Speech Commands/CIFAR-100 development TEST accuracy (\%).
Each uses the first 10 official classes, 128 TRAIN examples per class, common
features and 300 updates. Full 35/100-class heads are charged. This is not yet
full-dataset or class-incremental validation. Training wall times reflect the
original recorded possible concurrency, not isolated throughput.}
\label{tab:development-real}
\end{table}

The detached variant preserves the dynamic forward rule and parameter count,
but stops only the caller-input gradients from state and memory. CPU autograd
checks verify this intervention. It improves the poor speech optimization
trajectory, while not improving the image result; it is not a general solution.
Training and evaluation work, optimizer state, input caches, and failed
versions are recorded separately. The current Python-unrolled Block is slower
than the competitors despite smaller forward MAC proxies. Consequently these
observations do not establish a broad capability or complete-resource advantage.

\subsection{Questions and comparison conditions}

We distinguish three questions: whether new computations are acquired from
terminal feedback, whether later learning preserves their complete execution,
and whether topology independently changes the acquisition--retention--resource
tradeoff. Preservation includes input encoding, operation selection, composition,
and readout, rather than operation maps alone. All methods receive the same
permitted examples, terminal labels, stage order, and inference information.
Teacher transitions, intermediate execution states, correct routes, and
preconstructed task-specific primitives are excluded as free supervision.
Discovery, installation, observation-table passes, and protection are charged.

We separately compare matched information, matched resource budgets, and matched
attained quality. Equal update counts alone do not equate total learning work
when a method also performs program discovery or verification. Accuracy--parameter,
accuracy--training-cost, and accuracy--inference-cost frontiers expose these
differences. Changes to input adapters and output heads are recorded and charged;
a changed core learning rule constitutes a different architecture variant.

\subsection{Tasks and continual streams}

\begin{table}[t]
\centering
\small
\begin{tabular}{p{0.09\linewidth}p{0.40\linewidth}p{0.40\linewidth}}
\toprule
Layer & Task family & Principal measurement\\
\midrule
I & Six-stage streams at $N=8,16,32$ & Acquisition, stage-boundary and
within-stage retention, parameter growth\\
II & Noncommuting programs, long compositions, key--value memory & Order
sensitivity, unseen composition, delayed state use\\
III & Reversible, nonreversible, continuous, variable-length computation &
Generality of the same core Block\\
IV & Speech Commands v2 and Split CIFAR-100 & No-task-ID class-incremental
recognition and forgetting\\
V & Architecture competitors and component interventions & Competitive
frontiers and independent mechanism attribution\\
\bottomrule
\end{tabular}
\caption{Five required empirical layers. Success on one does not replace
measurements at another. The complete protocol is in Appendix~\ref{app:benchmarking}.}
\label{tab:main-evaluation-layers}
\end{table}

The finite computational streams expose observable initial states, token
programs, public task context, and terminal outputs. Context-dependent token
semantics make that public context available to every competitor. We disclose
short-word coverage, repeated examples, and every additional label access.
Both differentiable neural execution and native deployment are evaluated on
new and old tasks. Intermediate old-task minima and their subsequent recovery
are reported separately from stage-boundary retention.

The composition suite includes $AB\ne BA$, commutators, and unseen operation
combinations. Training lengths are at most eight; tests at lengths 16, 32, and 64
are reported separately. Key--value tasks require writing, overwriting,
reading, and delayed queries from terminal supervision. Generality tests cover
merging or erasing operations, continuous inputs/outputs, and variable-length
execution. Reversible lifting, extra workspace, or external memory counts
toward resources rather than making irreversibility free.

The following formal protocol remains prospective. Speech Commands v2 uses 35 classes in seven phases of five, the official
train/validation/test split, and common 16-kHz preprocessing. Split CIFAR-100
uses 100 fine classes in ten phases of ten, the official 50,000 training and
10,000 test images, and a fixed class-stratified 10\% training holdout for
validation. All methods share preprocessing and augmentation. Both tasks
recognize all seen classes without a task ID or class-group identifier at
inference. Their Memory-0 and Memory-500 conditions permit zero or at most 500
stored old examples in total. Stored features, teacher outputs, and statistics
are additionally disclosed. Main comparisons start from random initialization;
pretraining is supplemental with equal access across methods.

\subsection{Competitors and bounded development}

The primary competing family is the Transformer: vanilla, experience replay,
adapters, and replay plus adapters, at two prospectively selected scales.
Continual-learning comparisons include DEN, PackNet, replay, and ER-ACE or a
mature comparable class-incremental algorithm. Task-incremental PackNet is
separated from no-ID class-incremental rankings. Mamba-2 participates in
sequence/memory tests, Titans in applicable memory tasks, and HOPE/Nested
Learning where a verifiable implementation supports the interface. Ordinary
program learning, freezing, modular expansion, and constraint optimization
provide additional matched-information controls; they do not replace neural
architecture comparisons.

\begin{table}[t]
\centering
\small
\begin{tabular}{llll}
\toprule
Setting & Low budget & Medium budget & High budget\\
\midrule
Synthetic, updates/phase & B1:600 & B2:2,400 & B3:9,600\\
Real data, epochs/phase &10&50&200\\
\bottomrule
\end{tabular}
\caption{Prospective budget tiers available to every method. Discovery,
maintenance, and configuration-development work are additional reported costs.}
\label{tab:main-training-budgets}
\end{table}

Each main family has at most 24 development configurations. Development seed 0
selects candidates using frozen validation criteria; the top three are checked
with development seed 1. Configuration lock precedes final seeds 0, 1, 2, 3, 4.
Test and OOD results never select hyperparameters, retries, stopping points,
or configurations. Methods may use their appropriate optimizers and schedules.
Extra-budget trajectories remain separately labeled supplemental configurations.
Only localized, reproducible software, data, or training-logic defects qualify
as bug fixes; ordinary tuning is development.

Configuration validity depends on checked implementation, fidelity of any
claimed reproduction, information fairness, and compliance with the prospective
protocol, never on an accuracy threshold. High, intermediate, and random-level
scores are all valid outcomes under those conditions. A low-scoring adaptation
supports conclusions about that adaptation and budget, not an entire published
family. Unfaithful reproductions and untested implementations are explicitly
identified; they are not made invalid by low performance. Failed acquisition,
all failed seeds, and additional resource use are retained without extending
trials until a favorable result appears.

Learning curves distinguish failure to acquire a first task from forgetting
after later tasks. Training and validation diagnostics examine representation
limits, optimization, and credit assignment without using held-out tests to
choose a repair. A substantive architecture or optimizer revision receives a
new configuration identifier and preserves its predecessor's results. Formation
before gradient training, maintenance after training, and admission are separate
events, with their own label accesses and costs. This prevents installed
programs or confidence sharpening from being misreported as capabilities first
discovered by the gradient updates.

\subsection{Attribution, costs, and statistical analysis}

Same-architecture interventions remove reuse, sharing, state growth,
old-capability protection, graph constraints, and native compilation separately.
Topology-specific controls substitute ordinary constraints, disable cross-class
loop learning, disable within-class safe deformation, and compare ordinary
boundary optimization or discrete search. Information, certificates, and caches
are shared when applicable. Benefits from compression, replay, ordinary closure,
or confidence maintenance are not attributed to monodromy. A topology benefit
requires a measured residual effect under these concrete controls.

We record total, trainable, and frozen parameters; actual model, optimizer,
cache, replay, reference, and checkpoint storage; updates and sample accesses;
training FLOPs and GPU time; discovery/search and rejected proposals; growth,
maintenance, protection, verification, and reload. Controlled same-hardware
inference reports latency p50/p95, throughput, and peak memory. Concurrent wall
times are not ranked as isolated throughput, and overlapping timing/storage
fields are not summed twice. Parameters per newly retained capability include
the initial allocation and do not hide costs in nonparameter state.

Final accuracy, average incremental accuracy, forgetting, backward transfer,
acquisition cost, and equal-quality resource use are reported over all five
independent final seeds. We report means, standard deviations, 95\% confidence
intervals, and paired-seed differences; multiple primary tests receive
multiplicity correction. Resumed fragments of one lineage are not independent
replicates. Until these frozen-Primary measurements exist, this protocol
specifies the evaluation rather than asserting an experimental advantage.
